% =========================================================
% SECCIÓN 3.7 - ANÁLISIS ESTRUCTURAL Y CONSISTENCIA DEL BALANCE DE CARBONO
% =========================================================

Con el sistema de EDO formulado y los parámetros delimitados, esta sección verifica tres propiedades estructurales críticas para la validez del modelo: la consistencia dimensional de los flujos, el cierre del balance de masa de carbono y el comportamiento cualitativo del sistema en el dominio físico $\Omega$. El objetivo es demostrar que el modelo DEB escalado a reactor batch satisface las condiciones matemáticas mínimas que garantizan soluciones biológicamente interpretables, sin recurrir a las demostraciones formales de existencia, unicidad y estabilidad, que se remiten al Anexo \ref{anexo:formulacion_matematica}.

% ---------------------------------------------------------
\subsection{Homogeneidad dimensional de los flujos metabólicos}
\label{subsec:cap3_dimensional_resumen}
% ---------------------------------------------------------

La validez física de cualquier modelo cinético exige la homogeneidad dimensional en cada término de las ecuaciones de balance. Se verifica a continuación que todos los flujos del sistema comparten la unidad común de masa de carbono por unidad de tiempo (g C/min), garantizando que las cinco EDO sean internamente consistentes.

En la ecuación de balance del sustrato \eqref{eq:cap3_ode_sustrato}, la ingesta total $\dot{I}_{\mathrm{tot}}$ se construye como $N_{\mathrm{larvas}} \cdot \tilde{\imath}$, donde $N_{\mathrm{larvas}}$ es adimensional y $\tilde{\imath}$ tiene unidades de g C/larva/min. La forma funcional de Holling II \eqref{eq:cap3_holling} preserva esta unidad: $a_{\max}$ se expresa en g C$^{1/3}$/min, el factor de superficie $B^{2/3}$ en g C$^{2/3}$, y la funcional respuesta $S/(\tau_h + S)$ es adimensional, de modo que el producto completo resulta en g C/min. El factor de modulación térmica $f(T)$ de Sharpe--Schoolfield \eqref{eq:cap3_sharpe_schoolfield} es asimismo adimensional, al ser cociente de exponenciales con argumentos adimensionales $(T_A/T)$.

En el balance cuasi-estacionario de asimilados \eqref{eq:cap3_balance_asimilados}, cada término comparte la unidad g C/min: la asimilación $\dot{A} = \kappa_A \dot{I}_{\mathrm{tot}}$, el mantenimiento $\dot{M} = m B f(T)$ (min$^{-1} \cdot$ g C $\cdot$ adimensional), y los flujos de síntesis $\dot{G}_B, \dot{G}_L$ (g C/min). Los rendimientos $Y_B$ y $Y_L$ son adimensionales, de modo que $\dot{G}_B/Y_B$ y $\dot{G}_L/Y_L$ mantienen la homogeneidad dimensional. La prioridad $\kappa$ es adimensional, preservando la consistencia en la partición del excedente asimilado.

La ecuación de producción de \ch{CO2} \eqref{eq:cap3_ode_co2} acumula carbono disipado en g C/min. La conversión a flujo másico de \ch{CO2} en g/min mediante la ecuación \eqref{eq:cap3_conversion_masica} introduce el cociente de masas molares $M_{\ch{CO2}}/M_{\ch{C}} = 44.01/12.01 \approx 3.66$ (g \ch{CO2}/g C), factor adimensional que escala la tasa de carbono a la tasa de masa molecular de dióxido de carbono. Esta conversión es dimensionalmente consistente con la derivación del flujo $\dot{m}_{\ch{CO2}}(t)$ en el Capítulo \ref{chap:instrumentacion}, donde la pendiente de concentración $\mathrm{d}C/\mathrm{d}t$ (ppm/min) se transforma en flujo másico mediante la densidad molar del gas.

% ---------------------------------------------------------
\subsection{Conservación de masa de carbono}
\label{subsec:cap3_conservacion_resumen}
% ---------------------------------------------------------

El principio de conservación de la masa exige que la suma de las derivadas temporales de todos los compartimentos de carbono se anule cuando se consideran todas las fuentes y sumideros del sistema. Derivando la ecuación de balance global \eqref{eq:cap3_balance_global} respecto al tiempo, y despreciando la fracción no degradable $S_{\mathrm{no\_deg}}$ en el horizonte experimental de 15 días, se obtiene:
\begin{equation}
    \frac{\mathrm{d}S}{\mathrm{d}t} + \frac{\mathrm{d}E}{\mathrm{d}t} + \frac{\mathrm{d}B}{\mathrm{d}t} + \frac{\mathrm{d}L}{\mathrm{d}t} + \frac{\mathrm{d}C_{\ch{CO2}}}{\mathrm{d}t} = 0.
    \label{eq:cap3_conservacion_derivada}
\end{equation}

Sustituyendo las EDO \eqref{eq:cap3_ode_sustrato}--\eqref{eq:cap3_ode_co2} en la condición \eqref{eq:cap3_conservacion_derivada}, y utilizando el estado cuasi-estacionario de la reserva ($\mathrm{d}E/\mathrm{d}t \approx 0$), se obtiene una identidad algebraica que se verifica mediante manipulación directa de los flujos. En efecto, la ingesta $\dot{I}_{\mathrm{tot}}$ alimenta la asimilación $\dot{A} = \kappa_A \dot{I}_{\mathrm{tot}}$, la cual se distribuye en mantenimiento $\dot{M}$ y crecimiento $\dot{G}_B + \dot{G}_L$. La disipación total como \ch{CO2} es $\dot{M} + (1-Y_B)\dot{G}_B + (1-Y_L)\dot{G}_L$, de modo que la suma de todos los flujos de salida del sistema (crecimiento estructural, acumulación lipídica y \ch{CO2}) iguala exactamente la ingesta neta, cerrando el balance.

Es importante notar que el \ch{CH4} no aparece en este balance. Dado que el modelo asume metabolismo aeróbico exclusivo, cualquier carbono emitido como metano proviene de rutas fermentativas microbianas no incluidas en las EDO larval. En el Capítulo \ref{chap:metricas}, la detección de \ch{CH4} se utiliza para cuantificar la magnitud de esta vía anaeróbica paralela, permitiendo calcular la discrepancia acumulada $\Delta_{\mathrm{Total}}$ entre el balance teórico y la medición sensorial. La conservación de masa del modelo DEB, por tanto, se entiende como un balance aeróbico que se completa con la contribución diagnóstica del \ch{CH4} en la fase de evaluación ambiental.

% ---------------------------------------------------------
\subsection{Comportamiento cualitativo del sistema}
\label{subsec:cap3_comportamiento}
% ---------------------------------------------------------

El análisis cualitativo del sistema de EDO se centra en tres aspectos: la no negatividad de las variables de estado, la monotonicidad del sustrato y la acotación superior de la biomasa.

\paragraph{No negatividad.} El dominio $\Omega$ definido en la ecuación \eqref{eq:cap3_dominio} es invariante bajo el flujo del sistema siempre que las condiciones iniciales sean no negativas y los parámetros satisfagan las restricciones \eqref{eq:cap3_rest_kappaA}--\eqref{eq:cap3_rest_rhoconv}. En particular, cuando $S \to 0$, la funcional Holling II se anula ($\tilde{\imath} \to 0$), de modo que $\mathrm{d}S/\mathrm{d}t = 0$ y el sustrato no puede volverse negativo. Análogamente, la biomasa estructural $B(t)$ está protegida de la catabolismo en régimen de alimentación suficiente ($\dot{A} > \dot{M}$), y los lípidos $L(t)$ permanecen no negativos dado que $\dot{G}_L \geq 0$ en las mismas condiciones. Las demostraciones rigurosas de invarianza del dominio positivo se presentan en el Anexo \ref{anexo:formulacion_matematica}.

\paragraph{Monotonicidad del sustrato.} La ecuación \eqref{eq:cap3_ode_sustrato} garantiza que $\mathrm{d}S/\mathrm{d}t \leq 0$ para todo $t \geq 0$, dado que $\tilde{\imath} \geq 0$ por construcción. Por tanto, $S(t)$ es una función monótona decreciente que evoluciona desde $S_0$ hacia un valor residual $S_{\infty} \geq 0$. En la práctica experimental, el sustrato se renueva cada 48 h mediante adición de alimento fresco, de modo que la dinámica real presenta discontinuidades periódicas que el modelo promedia en una curva suavizada. Esta aproximación es válida cuando el período de renovación es corto respecto a la escala de crecimiento larval (15 días).

\paragraph{Acotación de la biomasa.} La biomasa estructural $B(t)$ no puede crecer indefinidamente: está limitada por el carbono total disponible $C_{\mathrm{tot}}$ y por la saturación de la ingesta cuando $S \gg \tau_h$. En el límite asintótico, cuando el sustrato se agota ($S \to 0$), la ingesta se anula y el crecimiento cesa ($\dot{G}_B \to 0$). La biomasa máxima teórica $B_{\max}$ se puede estimar a partir del balance de carbono: si toda la fracción asimilable del sustrato se convirtiera en biomasa estructural con rendimiento $Y_B$, se tendría $B_{\max} \approx \kappa_A Y_B S_0$, valor que en el reactor experimental de \SI{20}{\litre} con \SI{14}{\kilo\gram} de sustrato resulta del orden de \SI{3.5}{}--\SI{4.5}{\kilo\gram} C, consistente con los rendimientos de bioconversión reportados en la literatura \parencite{dienerConversionOrganicMaterial2011,parodiBioconversionEfficienciesGreenhouse2020a}.

\paragraph{Estabilidad del estado cuasi-estacionario.} La hipótesis de estado cuasi-estacionario para la reserva de asimilados ($\mathrm{d}E/\mathrm{d}t \approx 0$) es válida cuando la escala temporal de la dinámica de $E$ es mucho menor que la de $B$ y $L$. Dado que la reserva opera como un compartimento de amortiguamiento de alta velocidad —con constante de tiempo del orden de horas, frente a días para el crecimiento somático—, la separación de escalas está garantizada en el régimen experimental. El Anexo \ref{anexo:formulacion_matematica} presenta el análisis de estabilidad linealizada que justifica formalmente esta aproximación.

% ---------------------------------------------------------
\subsection{Confrontación con los datos del Capítulo \ref{chap:instrumentacion}}
\label{subsec:cap3_confrontacion}
% ---------------------------------------------------------

El comportamiento cualitativo predicho por el modelo DEB debe ser consistente con las observaciones empíricas del Capítulo \ref{chap:instrumentacion}. En particular, el flujo másico de \ch{CO2} medido por los sensores NDIR —expresado como $\dot{m}_{\ch{CO2}}(t)$ en g/min mediante la ecuación de conversión del Capítulo \ref{chap:instrumentacion}— debe correlacionarse con la disipación metabólica $\dot{D}(t)$ del modelo.

Bajo las condiciones del protocolo experimental (\SI{30}{\celsius} $\pm$ \SI{2}{\celsius}, HR $60\%$--$80\%$), el factor térmico $f(T)$ se mantiene cercano a su óptimo ($f \approx 0.9$--$1.0$), de modo que la tasa de mantenimiento $\dot{M}$ y la ingesta $\dot{I}$ evolucionan predominantemente según la biomasa estructural $B(t)$. Dado que $B(t)$ crece sigmoidalmente durante el ciclo larval, el modelo predice que $\dot{m}_{\ch{CO2}}(t)$ presentará un perfil unimodal: ascenso inicial acelerado por el crecimiento exponencial de la biomasa, seguido de un pico alrededor del estadio L5 y un descenso gradual conforme la población se aproxima a la prepupa y reduce su actividad metabólica. Este perfil es cualitativamente consistente con las series temporales de concentración $C(t)$ registradas en las ventanas cerradas de \SI{30}{\minute}, donde las pendientes $\Delta C/\Delta t$ alcanzan valores máximos entre los días 7 y 10 del ciclo \parencite{rossiEstimatingDynamicsGreenhouse2024a}.

La discrepancia entre el modelo teórico y la medición sensorial se define como:
\begin{equation}
    \delta(t) = \dot{m}_{\ch{CO2}}^{\mathrm{sensor}}(t) - \dot{m}_{\ch{CO2}}^{\mathrm{DEB}}(t),
    \label{eq:cap3_discrepancia}
\end{equation}
donde $\dot{m}_{\ch{CO2}}^{\mathrm{DEB}}(t)$ se calcula mediante las ecuaciones \eqref{eq:cap3_ode_co2} y \eqref{eq:cap3_conversion_masica}. En condiciones de aerobiosis estricta y homogeneidad del sustrato, se espera que $\delta(t)$ se mantenga dentro de la banda de incertidumbre instrumental del sensor ($\pm 12\%$). Desviaciones sistemáticas positivas ($\delta > 0$) sugieren una contribución microbiana elevada ($\dot{D}_{\mu}$ dominante), mientras que desviaciones negativas persistentes pueden indicar inhibición metabólica larval o formación de microambientes anaeróbicos no capturados por el modelo. El tratamiento estadístico de $\delta(t)$ y su vinculación con la detección de \ch{CH4} se desarrolla en el Capítulo \ref{chap:metricas}.

En resumen, esta sección ha verificado que el modelo DEB escalado a reactor batch satisface las condiciones estructurales de consistencia dimensional, conservación de masa y comportamiento cualitativo biológicamente plausible. Las demostraciones formales de existencia, unicidad y estabilidad de las soluciones, así como el análisis de sensibilidad de los parámetros, se remiten al Anexo \ref{anexo:formulacion_matematica}. El capítulo concluye en la siguiente sección con una síntesis de los hallazgos y el puente hacia la integración híbrida del Capítulo \ref{chap:hibrido}.
